\section{Problem formulation}\label{sec:problem}

\subsection{Hidden attractors}\label{sec:hidden}

Consider an autonomous system in $\mathbb{R}^n$: a flow $\dot x=F(x)$ or a map $x_{k+1}=F(x_k)$ with a smooth (or piecewise-smooth) right-hand side. Let $A$ be an attractor, $W(A)$ its basin, and $\mathcal{E}$ the set of equilibria (fixed points). The attractor $A$ is \emph{self-excited} if there exists $e\in\mathcal{E}$ such that $W(A)\cap U_\varepsilon(e)\neq\varnothing$ for every $\varepsilon>0$; otherwise, $A$ is \emph{hidden}.\cite{LeonovKuznetsov2013} Three consequences matter for the search. Trajectories launched from neighborhoods of equilibria do not reach a hidden attractor. In systems without equilibria and in systems with a single stable equilibrium, all nontrivial attractors are hidden. A hidden attractor may coexist with self-excited ones, and then its basin is an ``island'' that the unstable manifolds of the equilibria do not reach.

A continuous characteristic complements the binary definition. Let $\nu_\varepsilon$ be the uniform distribution on the union of the $\varepsilon$-neighborhoods of the equilibria (in an extended version, also on the $\varepsilon$-tubes around their unstable manifolds). We call
\begin{equation}
h_\varepsilon(A)=\Pr_{x\sim\nu_\varepsilon}\bigl[x\in W(A)\bigr]
\label{eq:heps}
\end{equation}
the \emph{hiddenness measure}. For a hidden attractor and a finite set of equilibria, $h_\varepsilon(A)=0$ for all sufficiently small $\varepsilon$, and the function $\varepsilon\mapsto h_\varepsilon$ shows the ``depth'' of hiddenness; for a continuum of equilibria, a common $\varepsilon$ is not guaranteed, and for systems without equilibria we set $h_\varepsilon=0$. A measured $h_\varepsilon=0$ is a necessary condition of hiddenness, not a sufficient one: it bounds the measure of $W(A)\cap U_\varepsilon(\mathcal E)$ from above by the sampling resolution. Below, it is measured for the benchmark systems (Sec.~\ref{sec:protocol}).

\subsection{Search as rare-event detection}\label{sec:rare}

Let $B\subset\mathbb{R}^n$ be a compact region of initial conditions with normalized Lebesgue measure $\mu$, $\mu(B)=1$. The target attractor $A^*$ has a basin $W^*$ of measure $p=\mu(W^*\cap B)\ll1$. The unit of cost is one integration of standard length $T=T_{\rm tr}+T_{\rm obs}$ (a transient segment and an observation window); for maps, it is one orbit. Below, one integration of one initial condition is also called one \emph{evaluation}. An algorithm $\mathcal{A}$ chooses initial points $x_1,x_2,\dots\in B$ adaptively: the point $x_k$ may depend on the results of the previous integrations. Let $\tau=\min\{k: x_k\in W^*\}$ be the index of the first hit. The quality of the algorithm is described by the detection probability and by the cost at confidence level $1-\delta$:
\begin{equation}
\begin{split}
&P_{\mathcal A}(N)=\Pr[\tau\le N],\\
&N_{1-\delta}(\mathcal A)=\min\{N: P_{\mathcal A}(N)\ge 1-\delta\}.
\end{split}
\label{eq:pn}
\end{equation}
The probability is taken over the internal randomness of the algorithm; throughout, $\delta=0.1$. The reference point is uniform independent sampling:
\begin{equation}
\begin{split}
&P_0(N)=1-(1-p)^N,\\
&N^0_{1-\delta}=\left\lceil\frac{\ln\delta}{\ln(1-p)}\right\rceil\approx\frac{\ln(1/\delta)}{p},\qquad N^0_{0.9}\approx\frac{2.3}{p}.
\end{split}
\label{eq:uniform}
\end{equation}
The gain factor $G_{1-\delta}(\mathcal A)=N^0_{1-\delta}/N_{1-\delta}(\mathcal A)$ is equivalent to an ``effective basin measure'' $p_{\rm eff}=pG$. On families with a controlled basin measure, we measure the scaling exponent $s$ in $N_{0.9}\propto p^{-s}$; for uniform sampling, $s=1$. A method has an \emph{order-of-growth gain} if $s<1$ with statistical significance and a \emph{constant-factor gain} if $s\approx1$ but $G>1$. A hit counts only if the candidate passes the identification and verification protocol of Sec.~\ref{sec:protocol}; transient chaos mistaken for an attractor is a failure of the method.

\subsection{Where a gain can come from}\label{sec:channels}

If an algorithm receives only the target indicator of each point, adaptivity cannot help in the worst case. Partition $B$ into $\lfloor1/p\rfloor$ cells of measure $p$ and let $W^*$ be a cell chosen uniformly at random. An algorithm that learns only whether $x_k\in W^*$ excludes at most one cell per query, so $\Pr[\tau\le N]\le Np$ and $N_{1-\delta}\ge(1-\delta)/p$ for every algorithm; uniform sampling is therefore within the factor $\ln(1/\delta)/(1-\delta)\approx2.6$ of the optimum. A gain is possible only through the information carried by the trajectories that did \emph{not} land in $W^*$. Among the sources of such information that cost one integration per point, we distinguish three channels (information about the vector field that needs no integration, such as harmonic balance or Lyapunov functions, is outside this setting).
\begin{itemize}
\item \textbf{Channel A: basin geometry.} Labels are spatially correlated. If the target basin adjoins already found boundaries between known basins, sampling near these boundaries is enriched with the target.
\item \textbf{Channel B: transients.} Trajectories converging to known attractors or escaping to infinity carry information about proximity to unseen invariant sets: the escape time, delays near saddles, and the mismatch between the early segment of a trajectory and its asymptotic regime. Near a basin boundary, a trajectory follows for a long time the stable manifold of the invariant set on the boundary (the edge state\cite{Skufca2006Edge,Schneider2007Edge,Vollmer2009Edge}), and its transient lengthens. Invariant sets on basin boundaries are classically located by straddle (bisection) methods that exploit exactly this effect,\cite{NusseYorke1989} and long transients are the signal of the stagger-and-step method for finding chaotic saddles.\cite{SweetNusseYorke2001}
\item \textbf{Channel C: prior sets.} Classical heuristics (nullclines, perpetual points, connecting curves) define a distribution of starting points that may be enriched with the target.
\end{itemize}
The proposed method (Sec.~\ref{sec:method}) uses all three channels. The cost model (Sec.~\ref{sec:costmodel}) indicates which channel works for which basin properties and provides an observable diagnostic.
